% GENERATED FILE. Edit canonical lessons, then run npm run generate:books.
% canonical-source-hash: fnv1a64:7f4a209d74622831
% canonical-lessons: ZH-C183-chuqu, ZH-C183-jinlai, ZH-C183-xiaban, ZH-C183-suomen, ZH-C183-tang

\chapter{In and Out}
\label{ch:zh-in-and-out}
\hlchaptermodality{183}

\begin{chapteropening}
\textbf{By the end of this chapter:} \emph{I can talk about going out, coming in, finishing work, locking up and lying down.}

Complete the last lesson of chapter 183: I can talk about going out, coming in, finishing work, locking up and lying down.
\end{chapteropening}

\section[\texorpdfstring{\zh{出去} (chūqù)}{出去 (chūqù)}]{\zh{出去} (chūqù) — to go outside; to go out of here}

\label{lesson:ZH-C183-chuqu}



\begin{quote}
Before the new one: say the Mandarin for home cooking, then the Mandarin for a boxed meal.
\end{quote}

\subsection*{You'll want to know: \zh{出去}}
\textbf{\zh{出去}} — \emph{chūqù} — \textquotedblleft{}to go outside; to go out of here\textquotedblright{}.

\begin{quote}
\textbf{\zh{外边很冷你还要出去吗}} \emph{wàibian hěn lěng nǐ hái yào chūqù ma} — It's cold out: are you still going out?
\end{quote}

\begin{cousinweb}
\textbf{\zh{出去}} (\emph{chūqù}) is built from \zh{出} 'to come out' and \zh{去} 'to go', both inherited from Middle Chinese, \emph{tsyhwit} and \emph{khjoH} (Guangyun, in Baxter's transcription): out and away from here.
\end{cousinweb}

\begin{grammarlens}[title={What you've built}]
One more word for saying what I do.
\end{grammarlens}

\subsection*{Your turn}
\begin{itemize}
\raggedright
  \item \emph{Say it:} \emph{chūqù}
  \item \emph{Say it:} \emph{chūqù}, once more
  \item \emph{Say it:} \emph{wàibian hěn lěng nǐ hái yào chūqù ma} — It's cold out: are you still going out?
  \item \emph{Recall:} say the Mandarin for home cooking, then the Mandarin for a boxed meal, then say \emph{chūqù} again
\end{itemize}

\begin{tcolorbox}[breakable,colback=teal!4,colframe=teal!35!black,title={Before you move on}]
What does \zh{出去} mean? (\textquotedblleft{}To go outside; to go out of here\textquotedblright{}.) Say it once more.
\end{tcolorbox}
\section[\texorpdfstring{\zh{进来} (jìnlái)}{进来 (jìnlái)}]{\zh{进来} (jìnlái) — to come inside; to come in}

\label{lesson:ZH-C183-jinlai}



\begin{quote}
Before the new one: say the Mandarin for a boxed meal, then the Mandarin for to go out of here.
\end{quote}

\subsection*{You'll want to know: \zh{进来}}
\textbf{\zh{进来}} — \emph{jìnlái} — \textquotedblleft{}to come inside; to come in\textquotedblright{}.

\begin{quote}
\textbf{\zh{请进来你的房间在这儿}} \emph{qǐng jìnlái nǐ de fángjiān zài zhèr} — Come on in: your room is here.
\end{quote}

\begin{cousinweb}
\textbf{\zh{进来}} (\emph{jìnlái}) is built from \zh{进} 'to go in' and \zh{来} 'to come', both inherited from Middle Chinese, \emph{tsinH} and \emph{loj} (Guangyun, in Baxter's transcription): in and towards here.
\end{cousinweb}

\begin{grammarlens}[title={What you've built}]
One more word for saying what I do.
\end{grammarlens}

\subsection*{Your turn}
\begin{itemize}
\raggedright
  \item \emph{Say it:} \emph{jìnlái}
  \item \emph{Say it:} \emph{jìnlái}, once more
  \item \emph{Say it:} \emph{qǐng jìnlái nǐ de fángjiān zài zhèr} — Come on in: your room is here.
  \item \emph{Recall:} say the Mandarin for a boxed meal, then the Mandarin for to go out of here, then say \emph{jìnlái} again
\end{itemize}

\begin{tcolorbox}[breakable,colback=teal!4,colframe=teal!35!black,title={Before you move on}]
What does \zh{进来} mean? (\textquotedblleft{}To come inside; to come in\textquotedblright{}.) Say it once more.
\end{tcolorbox}
\section[xiàbān]{xiàbān — to finish work; to come off a shift}

\label{lesson:ZH-C183-xiaban}



\begin{quote}
Before the new one: say the Mandarin for to go out of here, then the Mandarin for to come inside.
\end{quote}

\subsection*{You'll want to know: xiàbān}
\emph{xiàbān} — \textquotedblleft{}to finish work; to come off a shift\textquotedblright{}.

\begin{quote}
\emph{nǐ jǐ diǎn xiàbān} — What time do you finish work?
\end{quote}

Its characters come later; for now it is a word for the ear.

\begin{cousinweb}
\emph{xiàbān} is built from \zh{下} 'to go down' and \emph{bān}, both inherited from Middle Chinese, \emph{haeX} and \emph{paen} (Guangyun, in Baxter's transcription), which glosses the second as dividing tokens of jade; the sense 'a shift at work' is not given there.
\end{cousinweb}

\begin{grammarlens}[title={What you've built}]
One more word for saying what I do.
\end{grammarlens}

\subsection*{Your turn}
\begin{itemize}
\raggedright
  \item \emph{Say it:} \emph{xiàbān}
  \item \emph{Say it:} \emph{xiàbān}, once more
  \item \emph{Say it:} \emph{nǐ jǐ diǎn xiàbān} — What time do you finish work?
  \item \emph{Recall:} say the Mandarin for to go out of here, then the Mandarin for to come inside, then say \emph{xiàbān} again
\end{itemize}

\begin{tcolorbox}[breakable,colback=teal!4,colframe=teal!35!black,title={Before you move on}]
What does xiàbān mean? (\textquotedblleft{}To finish work; to come off a shift\textquotedblright{}.) Say it once more.
\end{tcolorbox}
\section[suǒmén]{suǒmén — to lock the door; to lock up}

\label{lesson:ZH-C183-suomen}



\begin{quote}
Before the new one: say the Mandarin for to come inside, then the Mandarin for to finish work.
\end{quote}

\subsection*{You'll want to know: suǒmén}
\emph{suǒmén} — \textquotedblleft{}to lock the door; to lock up\textquotedblright{}.

\begin{quote}
\emph{chūqù yǐqián bié wàng le suǒmén} — Don't forget to lock the door before you go out.
\end{quote}

Its characters come later; for now it is a word for the ear.

\begin{cousinweb}
\emph{suǒmén} is built from \emph{suǒ} 'a lock' and \zh{门} 'door', both inherited from Middle Chinese, \emph{swaX} and \emph{mwon} (Guangyun, in Baxter's transcription), which glosses the first 'an iron lock'.
\end{cousinweb}

\begin{grammarlens}[title={What you've built}]
One more word for saying what I do.
\end{grammarlens}

\subsection*{Your turn}
\begin{itemize}
\raggedright
  \item \emph{Say it:} \emph{suǒmén}
  \item \emph{Say it:} \emph{suǒmén}, once more
  \item \emph{Say it:} \emph{chūqù yǐqián bié wàng le suǒmén} — Don't forget to lock the door before you go out.
  \item \emph{Recall:} say the Mandarin for to come inside, then the Mandarin for to finish work, then say \emph{suǒmén} again
\end{itemize}

\begin{tcolorbox}[breakable,colback=teal!4,colframe=teal!35!black,title={Before you move on}]
What does suǒmén mean? (\textquotedblleft{}To lock the door; to lock up\textquotedblright{}.) Say it once more.
\end{tcolorbox}
\section[tǎng]{tǎng — to lie down}

\label{lesson:ZH-C183-tang}



\begin{quote}
Before the new one: say the Mandarin for to finish work, then the Mandarin for to lock up.
\end{quote}

\subsection*{You'll want to know: tǎng}
\emph{tǎng} — \textquotedblleft{}to lie down\textquotedblright{}.

\begin{quote}
\emph{wǒ hěn lèi xiǎng tǎng yíhuìr} — I'm tired; I want to lie down for a bit.
\end{quote}

Its characters come later; for now it is a word for the ear.

\begin{cousinweb}
\emph{tǎng} No dictionary we use records an older form, so its history is not settled.
\end{cousinweb}

\begin{grammarlens}[title={What you've built}]
One more word for saying what I do.
\end{grammarlens}

\subsection*{Your turn}
\begin{itemize}
\raggedright
  \item \emph{Say it:} \emph{tǎng}
  \item \emph{Say it:} \emph{tǎng}, once more
  \item \emph{Say it:} \emph{wǒ hěn lèi xiǎng tǎng yíhuìr} — I'm tired; I want to lie down for a bit.
  \item \emph{Recall:} say the Mandarin for to finish work, then the Mandarin for to lock up, then say \emph{tǎng} again
\end{itemize}

\begin{tcolorbox}[breakable,colback=teal!4,colframe=teal!35!black,title={Before you move on}]
What does tǎng mean? (\textquotedblleft{}To lie down\textquotedblright{}.) Say it once more.
\end{tcolorbox}
